%=====================================================================
\clearpage
\hdr{J4}{光栅常量、线数与角色散}

\subhdr{① 这一环交出什么}

\begin{center}
\begin{minipage}{\linewidth}
\begin{tabular}{@{}p{3.6cm}p{4.2cm}p{6.2cm}@{}}
\toprule
\thead{产出} & \thead{合格判据} & \thead{数据来源} \\
\midrule
光栅常量 $d$、线数 $N$、角色散 $D$ & $N$ 接近整数（$100/300/600$）；一级 $D\approx1'/\mathrm{nm}$ & 全部用 J3 的 $\varphi_k$，\textbf{不需要新测量} \\
\bottomrule
\end{tabular}
\end{minipage}
\end{center}

用时约 30 分钟，全部是计算与核对。这一环可以离开仪器做，但\textbf{如果发现 \S ⑤ 的自检不过，趁仪器还在赶紧回去补测}。

\subhdr{② 为什么这样能算}

\textbf{光栅常量}：光栅方程里已知 $\lambda$、测得 $\varphi_1$，直接解出

\[
d=\frac{\lambda}{\sin\varphi_1}.
\]

\textbf{线数}：出厂标的是「每毫米多少条缝」$N$，而公式里用的是缝距 $d$。两者的关系是 $d=1/N$（$d$ 以 mm 计），也就是

\[
N=\frac{1}{d}\quad(d\ \text{以 mm 计}),\qquad
N=\frac{1}{1000\,d}\quad(d\ \text{以 m 计}).
\]

\textbf{角色散}：定义为「单位波长差对应的角间距」，

\[
D=\frac{\Delta\varphi}{\Delta\lambda}
\quad\text{（直接测量法：用一级的紫线与绿线）},
\]
\[
D=\frac{\dd\varphi}{\dd\lambda}=\frac{k}{d\cos\varphi}
\quad\text{（由光栅方程求导：$d\sin\varphi=k\lambda$ 两边对 $\lambda$ 求导）}.
\]

两个式子应互相印证。$300$ 线/mm 的光栅，一级角色散约 $1'/\mathrm{nm}$——记住这个量级，它是判断对错的锚。

\subhdr{③ 不需要新操作}

这一环的全部输入是 J3 的八个方位读数。唯一要确认的是：\textbf{八个读数对应同一块光栅、同一套调节}；中途若动过 16 号或重新调过 $b_1b_2$，前后数据不能混着用。

\subhdr{④ 汇总表（把 J3 的结果抄进来算）}

\begin{center}
\begin{minipage}{\linewidth}
\begin{tabular}{@{}p{3.4cm}p{3.4cm}p{3.4cm}p{4.4cm}@{}}
\toprule
\thead{量} & \thead{公式} & \thead{测得值} & \thead{与预期是否相符} \\
\midrule
\rowstrut $d$（一级绿线） & $\lambda/\sin\varphi_1$ & & $3333$ nm（$300$ 线/mm） \\
\rowstrut $d'$（二级绿线核对） & $2\lambda/\sin\varphi_2$ & & 与 $d$ 差 $<1\%$ \\
\rowstrut $N$ & $1/d$ & & $100/300/600$ \\
\rowstrut $D$（一级紫绿） & $\Delta\varphi/\Delta\lambda$ & & $\approx1.0\ '/\mathrm{nm}$ \\
\rowstrut $D$（解析式核对） & $k/(d\cos\varphi)$ & & 与上一行差 $<1\%$ \\
\bottomrule
\end{tabular}
\end{minipage}
\end{center}

\subhdr{⑤ 当场算：一步一步}

\textbf{第 1 步：$d$ 用一级绿线。}手册明确指定一级绿线（$\lambda=546.1$ nm）：

\[
d=\frac{546.1\ \mathrm{nm}}{\sin 9^{\circ}26'}=\frac{546.1}{0.16387}=3333\ \mathrm{nm}=3.333\ \mu\mathrm{m}.
\]

\textbf{第 2 步：用二级做独立核对。}

\[
d'=\frac{2\lambda}{\sin\varphi_2}=\frac{1092.2\ \mathrm{nm}}{\sin 19^{\circ}08'}=\frac{1092.2}{0.3277}=3333\ \mathrm{nm}\ \checkmark
\]

理论上 $d'=d$。实测应相差在 $1\%$ 以内；差很多说明某一级角度测偏了（多半是栅面没垂直光轴，或谱线不等高）。

\textbf{第 3 步：换成线数并与出厂值对照。}

\[
N=\frac{1}{d}=\frac{1}{3.333\times10^{-3}\ \mathrm{mm}}=300\ \text{线/mm}.
\]

$N$ 应接近整数。算出 $293$ 或 $307$ 是正常的测量偏差；算出 $2930$ 说明单位换算错了 10 倍。

\textbf{第 4 步：角色散用一级紫绿两线。}

\[
D=\frac{\varphi_{1,\text{绿}}-\varphi_{1,\text{紫}}}{546.1-435.8}
=\frac{9^{\circ}26'-7^{\circ}31'}{110.3\ \mathrm{nm}}
=\frac{115'}{110.3\ \mathrm{nm}}=1.04\ '/\mathrm{nm}.
\]

与解析式 $D=\dfrac{1}{d\cos 9^{\circ}26'}=\dfrac{1}{3333\times0.9865}=1.045'/\mathrm{nm}$ 一致 $\checkmark$。

\warnbox{\textbf{三处最容易出错的地方}：① 把 $N$（线/mm）当成 $d$ 直接代进 $d\sin\varphi=k\lambda$；② $d$ 用 $\mu$m、$\lambda$ 用 nm，忘了换算（$1\ \mu\mathrm{m}=1000$ nm）；③ 漏掉 $N=1/(1000d)$ 里的 $1000$，算出的 $d$ 差一千倍。统一做法：\textbf{先把 $d$ 换到 nm}，再与 $\lambda$ 一起代。}

\begin{center}
\begin{minipage}{\linewidth}
\begin{tabular}{@{}p{2.2cm}p{3.0cm}p{3.4cm}p{4.6cm}@{}}
\toprule
\thead{$N$（线/mm）} & \thead{$d=1/N$} & \thead{一级绿线 $\varphi_1$} & \thead{量级感觉} \\
\midrule
$100$ & $10.000\ \mu\mathrm{m}$ & $3.13^{\circ}$ & 一级很近，容易与零级混 \\
$300$ & $3.333\ \mu\mathrm{m}$ & $9.43^{\circ}$ & 本实验最常见 \\
$600$ & $1.667\ \mu\mathrm{m}$ & $19.13^{\circ}$ & 一级已经很散开 \\
\bottomrule
\end{tabular}
\end{minipage}
\end{center}

\textbf{光栅光谱与棱镜光谱的区别}（口试第 11 问）：

\begin{center}
\begin{minipage}{\linewidth}
\begin{tabular}{@{}p{2.4cm}p{5.4cm}p{5.4cm}@{}}
\toprule
\thead{比较项} & \thead{光栅} & \thead{棱镜} \\
\midrule
分光机制 & 多缝干涉（$d\sin\varphi=k\lambda$） & 折射率随波长变（色散） \\
谱线排列 & 长波偏折大（\textbf{红在外}），近似等间距 & 短波偏折大（\textbf{紫在外}），不等间距 \\
级次 & 有多级，高级次会重叠 & 只有一套光谱 \\
\bottomrule
\end{tabular}
\end{minipage}
\end{center}

\subhdr{⑥ 写进报告的一句话}

\callout{
「由一级绿线得光栅常量 $d=3.333\ \mu\mathrm{m}$，对应线数 $N=300$ 线/mm；用二级复核得 $d'=3.333\ \mu\mathrm{m}$，两者相差小于 $1\%$；由一级紫绿两线得角色散 $D=1.04'/\mathrm{nm}$，与解析式 $k/(d\cos\varphi)$ 一致。」
}
\warnbox{
\textbf{台前只看这一段}：$d=\lambda/\sin\varphi_1$（用一级\textbf{绿}线，$\lambda=546.1$ nm）；$d'=2\lambda/\sin\varphi_2$ 核对；$N=1/d$（$d$ 换 mm）；$D=\dfrac{\varphi_{\text{绿}}-\varphi_{\text{紫}}}{110.3\ \mathrm{nm}}$；先把 $d$ 换到 nm 再代。
}